\rightline{（文責：藻浦乗成）}
\section{ボードゲームをなんとよぶ}

　熊野寮は広大です。ひとくくりに寮生といっても、その趣味嗜好は多岐にわたり、全寮生に通底する指向性を示すことは困難です。しかし、ある程度全寮的に親しまれている余暇活動も少なからず存在します。ボードゲームもそのひとつです。筆者は幼少のころからボードゲームにふれ、寮生活のあいだにもさまざまなボードゲームを遊んできました。そのなかでボードゲームというジャンルの複雑さにふれ、「一般に想定される全てのボードゲームを支持しうる包括的な呼称は存在しうるか」という問題に直面してきました。結論を先取すると、持論としては「不可能」の三文字で終わるのですが、その道筋の試行錯誤において現れるボードゲームの多様さを伝えたいと思い、こうして書いている次第です。個人的な瑣末なこだわりではありますが、寮生活の一端にふれるものとしてご寛恕いただければ幸いです。

\subsection{ボードゲーム}
　最もスタンダードな呼称かと思います。本記事でもひとまずこの語を用います。しかしいちおう問題点を指摘するなら、いわゆるトランプゲームやダイスゲームのような、盤面上で行われないゲームの扱いについて解釈が分かれかねない点が挙げられます。よく知られたものでは、「ニムト」[1]「コヨーテ」[2]などもボードを使用せずカードのみによって進行します。特に混乱を招くのは、「カードゲーム」という呼称が「マジック：ザ・ギャザリング（M:tG）」[3]「遊☆戯☆王オフィシャルカードゲーム（OCG）」[4]などのゲームについて用いられていることです。「ニムト」[1]「コヨーテ」[2]は製作者・販売者・デザインなどから明らかにボードゲームの文脈のなかに位置付けられるものですが、「M:tG」[3]「OCG」[4]はトレーディングカードゲーム（TCG）として全く異なる文脈をもちます。

\subsection{アナログゲーム}
　電子機器を用いないゲームとの含意を込めて用いられ、類似したものとして「非電源ゲーム」との呼称もみられます。主にコンピューターゲーム・デジタルゲームと対比する意図が見てとれ、教育の文脈で多く用いられている印象があります。その起源にレトロニムとしての要素を含む以上、包括的な語としては欠陥を残します。第一に、そもそもの「アナログ」という語の意味と全く結びついていません。また、一部のウォーゲームや「くるりんパニック」[6]などのボードゲームは電池・電源を使用します。さらに「アンロック」[7]など、ボードゲームの構成やデザインを踏まえつつ、スマートフォンとの連携によって新しい体験を提供するゲームも現れています。

\subsection{テーブルゲーム}
　卓上で顔を突き合わせて遊ぶことを踏まえた呼称です。しかし、これも重箱の隅をつつくような指摘ではありますが、例えば「キャプテン・リノ：巨大版」[8]のような卓上で遊ぶことを想定していないゲームや、「いかさまゴキブリ」[9]のような机の外に重心が置かれたゲームもいくつか存在します。またこの語は、ボードゲームと一定の距離を持って発展したテーブルトーク・ロールプレイングゲーム（TRPG）の文脈を含むこともしばしばあり、一般的なボードゲームの指示範囲・系統と多少のずれを含みます。

\subsection{マインドスポーツ}
　スポーツとのアナロジーで知的能力を用いる競技をさす呼称です。「マインドスポーツオリンピアード」などの競技会が中心的な位置を占め、囲碁・チェス・ポーカーなどの古典的なゲームのほか「カタンの開拓者たち」[10]「テラフォーミング・マーズ」[11]など現代のボードゲームが扱われることもあります。身体競技に対する頭脳競技としての側面が強調されるため、「ディクシット」[12]のような感性に依拠するゲーム、「ボムバスターズ」[13]のような協力型ゲームを扱う蓋然性は非常に低いと言わざるを得ません。

\subsection{ドイツゲーム}
　現代のボードゲームがドイツを中心に展開したことに由来する呼称です。所要時間が30〜90分程度であり、途中での脱落がなく、戦略が求められ、かつ複雑すぎず、具体的なモチーフをもつ、といった特徴が挙げられます。1979年以来毎年ドイツで開催される「ドイツ年間ゲーム大賞（Spiel des Jahres）」が注目されたことでジャンルとして確立されました。1995年に「カタンの開拓者たち」が大賞を受賞して以来、より複雑なゲームも取り上げられるようになり、2011年からは「エキスパート賞（Kennerspiel des Jahres）」も併設されています。しかし、ボードゲーム市場の多極化、隣接分野との相互作用は、歴史的文脈をはるかに越える幅を持って進行しつつあります。

　以上のように「ボードゲーム」の呼称について概観すると、ボードゲームの多義性・多様性が注目されます。ここまで折にふれて挙げてきた「くるりんパニック」[6]「ディクシット」[12]「ボムバスターズ」[13]などのゲームも、それぞれ特異性をもちながら、先述した「ドイツ年間ゲーム大賞」において部門賞・大賞受賞の栄誉に浴しています。さらに、ギガミック社[14]やネスターゲームス[15]の製品に代表される具体的なモチーフをもたないゲーム（いわゆるアブストラクトゲーム[16]）、ドイツゲーム確立以前に制作された「モノポリー」[17]「スクラブル」[18]「アクワイア」[19]などのゲーム、「シェフィ」[20]のような一人用ゲームなど分類を困難にする要素は無数に存在します。このような状況下ではいずれの語を用いるにせよ、「××でない」「××に限る」といった排他的な定義をうちたてるのではなく、制作形態・流通経路・デザインなどあらゆる面から”ボードゲーム”に近しいものを博捜し拾い上げていく運用を実施するほうが有意義に思えます。

　……さて、改めて言い直すまでもなく、この結論は寮内連帯と相似形で組み立てられたものです。多様な人員を擁する熊野寮において、全寮生に通底する指向性を示すことは困難です。互いに規範・意向を押し付けずに済むよう、共有できる事象を出発点として合意の取れる結論を探り、多様な寮生を包摂できる弾力性のある自治を行うことが理想であると筆者は考えます。幸いというべきか、寮生の数は現在流通するボードゲームの種類よりはるかに少なく、われわれは「同じ建物に居住する」という明確な共通点を持ちます。そして少なくとも、こんな瑣末なことに偏屈なこだわりを持っている筆者が暮らしていけている程度には、熊野寮は懐の広い場所です。これから入寮するみなさんにも一定の寛容を求めると同時に、同様の寛容をもって皆さんを迎えることをここに約束します。

参考：高橋浩徳『ボードゲームの近現代史』「大阪商業大学アミューズメント産業研究所紀要」20, 2018

\begin{tcolorbox}[colback=white,colframe=black,title=本文中で挙げたゲームについての注記]
[1] Wolfgang Kramer, 1994 のろまな牛の行列を押し付け合うゲーム。計画性と危険予知が求められる。他プレイヤーの行動で予定が狂わされ、たいへんなことになる。\\\relax
[2] Spartaco Albertarelli, 2003 数字の書かれたカードを一枚ずつ持ち、自分以外のカードを見て自分のカードを推測するゲーム。簡単なルールながら深い推理とブラフが楽しめる。\\\relax
[3] Richard Garfield, 1993 世界初のTCG。リチャード・ガーフィールドはのちに「キング・オブ・トーキョー」（2011）等のボードゲームも作成している。\\\relax
[4] Konami Digital Entertainment, 1999 高橋和希の漫画『遊☆戯☆王』に登場するTCG「マジック・アンド・ウィザーズ」を再現したものとして制作・発売される。\\\relax
[5] ジャンルとしてのTCGは「プレイヤーがそれぞれ収集・構築したデッキを持つ」ことによって特徴づけられる。「ドミニオン」（Donald X. Vaccarino, 2008）はTCGの要素をボードゲームで再現した「デッキ構築ゲーム」として制作されたものである。\\\relax
[6] Carol Wiseley, 1992 タイミングよくペダルを叩いて、回転しながら襲ってくる空賊から農場のニワトリを守るゲーム。「くるりんパニック」やバランスゲーム系統など、器用さや素早さを用いて遊ぶゲームはしばしば「アクションゲーム」とよばれる。\\\relax
[7] Alice Carroll, Thomas Cauët, Cyril Demaegd, 2017 カードに描かれた道具や情報を駆使して状況を解決する。脱出ゲームや謎解き、ゲームブックの影響を強く受けている。\\\relax
[8] Scott Frisco \& Steven Strumpf, 2011 バランスを保って壁と屋根を積み上げ、家を建てるゲーム。巨大版は2015年発売、高さ3mにも達する家を作ることができる。\\\relax
[9] Emely Brand \& Lukas Brand, 2011 ルールに従ってカードを出し、手札をなくしていくゲーム。ただし、警備役にバレなければこっそりカードを隠したり捨てたりしてもよい。\\\relax
ゲームマスターがシナリオを管理し、プレイヤーたちがキャラクターを演じながら進行するゲーム。キャラクターの技能や自由度の高いロールプレイによって探索を進める。\\\relax
[10] Klaus Teuber, 1995 無人島を開拓して拠点を築き、拠点から得た資源を使って新たな拠点を築いていくゲーム。「拡大再生産ゲーム」とよばれるジャンルの代表的な作品。\\\relax
[11] Jacob Fryxelius, 2016 火星を開拓し居住可能な環境に近づけていく拡大再生産ゲーム。多様な変数を管理し多彩なカードを活用して勝利を狙う。いわゆる重量級ゲームの代表。\\\relax
[12] Jean-Louis Roubira, 2008 絵画を命名するゲーム。示唆的なイラストに、ほどよく曖昧で共感性がありつつあからさまでないタイトルをつける。\\\relax
[13] 林 尚志, 2025 プレイヤー全員で協力して正しく導線を切り、爆弾を解除するゲーム。日本人の作品では初のドイツ年間ゲーム大賞における大賞受賞作。\\\relax
[14] 戦後ドイツで発展したため戦争などのモチーフは忌避される傾向にあり、18世紀以来開発されてきた戦争シミュレーションゲーム（ウォーゲーム）とは一定の懸隔がある。\\\relax
[15] 「クアルト」（Blaise Muller, 1991）「コリドール」（Mirko Marchesi, 2007）などデザイン性に優れたクラシカルなゲームを多く販売するフランスのメーカー。\\\relax
[16] 「ヤバラス」（Cameron Browne, 2007）「ブリンク」（Néstor Romeral Andrés, 2020）などシンプルながら奥行きのあるゲームを多く販売するスペインのメーカー。\\\relax
[17] 原義としては囲碁やオセロのような「具体的な形を象らないゲーム」をさすが、現在は「二人対戦・公開情報のみ・運要素なし」の盤上ゲームという意味で用いられている。\\\relax
[18] Charles Darrow, 1935 双六の盤面上の土地を購入し、不動産から収入を得るゲーム。\\\relax
[19] James Brunot, 1948 クロスワードのように縦横に文字を並べて単語を作るゲーム。\\\relax
[20] Sid Sackson, 1962 ホテルの株主となり規模拡大・吸収合併を繰り返して儲けるゲーム。\\\relax
[21] ポーン, 2013 イベントカードを駆使して羊を1000匹まで増やしていくゲーム。
\end{tcolorbox}
